\documentclass[11pt,a4paper]{article}

\usepackage[utf8]{inputenc}
\usepackage[T1]{fontenc}
\usepackage[french]{babel}
\usepackage{lmodern}
\usepackage[a4paper,margin=2.5cm]{geometry}
\usepackage{microtype}
\usepackage{booktabs}
\usepackage{array}
\usepackage{longtable}
\usepackage{xcolor}
\usepackage{hyperref}
\usepackage{enumitem}
\usepackage{fancyhdr}
\usepackage{listings}
\usepackage{amsmath,amssymb}

\definecolor{teal}{RGB}{19,130,123}
\definecolor{navy}{RGB}{18,48,74}
\definecolor{lightgray}{RGB}{246,248,250}
\hypersetup{
  colorlinks=true,
  linkcolor=teal!60!black,
  urlcolor=teal!60!black,
  citecolor=teal!60!black,
  pdftitle={Classification du sentiment financier par entropie et compression},
  pdfauthor={Votre nom}
}
\pagestyle{fancy}
\fancyhf{}
\lhead{\small Sentiment financier, information et compression}
\rhead{\small Rapport de projet}
\cfoot{\thepage}
\setlength{\headheight}{14pt}
\setlength{\parskip}{0.45em}
\setlength{\parindent}{0pt}
\setlist[itemize]{leftmargin=1.5em,itemsep=0.2em,topsep=0.2em}
\setlist[enumerate]{leftmargin=1.7em,itemsep=0.2em,topsep=0.2em}
\lstset{
  basicstyle=\ttfamily\small,
  backgroundcolor=\color{lightgray},
  frame=single,
  rulecolor=\color{black!15},
  breaklines=true,
  columns=fullflexible,
  showstringspaces=false
}

\title{\vspace{-1.5cm}\color{navy}\Huge Sentiment financier par entropie et compression\\[0.4em]
\Large Une étude expérimentale inspirée de l'information algorithmique et du principe MDL}
\author{Votre nom\\[0.3em]\normalsize Votre établissement}
\date{\normalsize Rapport de projet}

\begin{document}
\maketitle
\thispagestyle{empty}

\begin{abstract}
Ce projet étudie une question pratique à la croisée du traitement automatique
des langues et de la théorie de l'information : dans quelle mesure peut-on
classer le sentiment de textes financiers à l'aide de représentations et de
modèles compacts, plutôt qu'en s'appuyant uniquement sur un grand modèle
pré-entraîné ? Deux pistes légères sont mises en œuvre et comparées à FinBERT :
un classifieur Bernoulli Naive Bayes dont les mots sont sélectionnés et
pondérés par leur gain d'information, et un classifieur qui rapproche la
classification du principe Minimum Description Length (MDL) au moyen de
dictionnaires zlib propres à chaque classe. Une chaîne d'évaluation commune
est appliquée à Financial PhraseBank : 3\,448 phrases dédupliquées, accord
d'annotation d'au moins 75\,\%, partition stratifiée 80/20 et graine fixée à
42. Sur ce seul test pilote, le modèle zlib atteint une accuracy de 72,4\,\%
et un F1 macro de 50,9\,\%, contre 93,9\,\% et 93,1\,\% pour FinBERT. Le
classifieur zlib est plus rapide et son état sauvegardé est beaucoup plus
petit, mais il détecte mal la classe négative. Ces résultats illustrent un
compromis entre qualité et coût ; ils ne démontrent ni équivalence des
approches ni supériorité générale de la compression.
\end{abstract}

\vfill
\textbf{Mots-clés :} apprentissage automatique, induction, entropie,
compression, MDL, sentiment financier, Financial PhraseBank, FinBERT.
\newpage
\tableofcontents
\newpage

\section{Résumé du travail}

\subsection{Question étudiée}
Les textes financiers sont courts, mais leur tonalité peut être difficile à
interpréter automatiquement. Une variation de vocabulaire, une négation ou
une expression spécifique au domaine peut inverser le sens d'une phrase.
Les modèles de langue pré-entraînés comme FinBERT sont capables de traiter
ces phénomènes avec une qualité élevée, mais ils ont un coût de calcul et une
taille supérieurs à ceux de classifieurs lexicaux compacts.

Le projet compare donc deux façons d'exploiter l'information contenue dans
les textes sans construire un grand modèle neuronal. La première mesure
l'association entre les mots et les étiquettes. La seconde compare les coûts
de description obtenus en compressant un texte avec des informations issues
des différentes classes. L'idée directrice est celle du cours :
une représentation utile conserve les régularités informatives et évite de
décrire à nouveau tout ce que le modèle peut déjà expliquer.

\subsection{Ce qui a été construit}
Le dépôt contient trois voies de classification, une procédure de comparaison
reproductible et ses artefacts :
\begin{itemize}
  \item \textbf{Entropie pondérée :} tokenisation, représentation binaire de
  présence/absence, sélection des mots par gain d'information et modèle
  Bernoulli Naive Bayes pondéré.
  \item \textbf{Compression zlib/MDL :} un dictionnaire est construit pour
  chaque classe à partir du train ; chaque phrase de test est évaluée en
  fonction du coût du dictionnaire, de son coût résiduel compressé et de
  l'a priori de classe.
  \item \textbf{Référence FinBERT :} modèle financier pré-entraîné évalué en
  inférence, sans ajustement sur le corpus du benchmark.
  \item \textbf{Évaluation :} baseline majorité, métriques par classe,
  matrice de confusion, durées, taille sérialisée ou taille des poids, et
  rapport JSON/Markdown contenant le protocole.
\end{itemize}

Le résultat central est nuancé : les méthodes compactes diminuent les coûts
mesurés, mais leur qualité demeure nettement inférieure à FinBERT sur
l'expérience disponible. La méthode zlib est la meilleure des deux méthodes
légères, mais son rappel des exemples négatifs est particulièrement faible.
Ce constat est essentiel : un modèle compact n'est utile que si la qualité
qu'il conserve est suffisante pour l'usage visé.

\section{Lien avec le cours : apprendre, c'est choisir une description}

\subsection{Induction à partir d'exemples}
L'induction consiste à inférer une règle générale à partir d'un nombre fini
d'observations. Pour la tâche étudiée, les observations sont des phrases
financières accompagnées d'une étiquette : \emph{positive}, \emph{negative}
ou \emph{neutral}. Le classifieur doit produire une règle qui généralise aux
phrases non vues.

Ce problème est sous-déterminé : de nombreuses règles peuvent expliquer un
échantillon d'entraînement. Une règle qui mémorise chaque phrase fonctionne
sur les exemples connus, mais ne permet pas nécessairement de classifier une
nouvelle phrase. À l'inverse, une règle trop simple peut ignorer les indices
linguistiques qui distinguent les classes. L'enjeu est donc de choisir une
représentation et un modèle qui équilibrent simplicité et adéquation aux
données.

Dans le projet, cet équilibre se manifeste concrètement dans la limite du
vocabulaire de la méthode entropique, dans la taille limitée des dictionnaires
zlib et dans la comparaison de leurs performances à un modèle pré-entraîné.
La règle de sélection n'est pas universelle : elle dépend du corpus, des
représentations permises et du coût retenu.

\subsection{Analogie et limites de la portée du projet}
Le chapitre du cours présente aussi l'analogie : une transformation apprise
sur un exemple peut être transférée à une structure similaire, comme une
règle morphologique appliquée à un nouveau mot. Le classifieur de sentiment
ne résout pas directement des analogies du type $A:B::C:D$ et ne met pas en
œuvre le programme spécialisé de génération de chaînes décrit dans le cours.
Le rapprochement est conceptuel : dans les deux cas, on cherche une régularité
qui explique plusieurs exemples plutôt que de les traiter isolément.

La distinction est importante pour interpréter le livrable. Le projet est une
application de l'induction et de la compression à la classification de
sentiment, non une reproduction de l'ensemble des exercices du chapitre. De
même, une anomalie pourrait en principe être associée à un coût résiduel élevé
ou à une faible confiance, mais aucun détecteur d'anomalies n'a été évalué
dans les résultats présentés ici.

\subsection{Principe MDL}
Le principe MDL choisit le modèle qui minimise la longueur totale de la
description du modèle et des observations qui ne sont pas expliquées par
celui-ci. Sous une forme schématique :
\[
  M^\star = \arg\min_M \left[C(M)+C(D\mid M)\right],
\]
où $C(M)$ désigne le coût de description du modèle et $C(D\mid M)$ celui des
données conditionnellement à ce modèle.

Un classifieur qui explique bien les régularités de ses exemples peut réduire
la quantité d'information qu'il reste à encoder. Mais ajouter une règle
complexe n'est intéressant que si le gain sur les résidus compense le coût de
cette règle. La version pratique de MDL utilisée ici s'inspire de ce
compromis sans prétendre calculer la complexité de Kolmogorov ni trouver un
code universel optimal. zlib fournit un code effectif et calculable, mais le
résultat dépend de ses choix d'implémentation et du protocole d'apprentissage.

\section{Données et protocole expérimental}

\subsection{Financial PhraseBank}
Le benchmark utilise la version de Financial PhraseBank correspondant à un
accord entre annotateurs d'au moins 75\,\%. Ce corpus contient des phrases
financières étiquetées en trois catégories. Les doublons textuels de même
étiquette sont dédupliqués par le chargeur ; un doublon avec une étiquette
contradictoire déclenche une erreur plutôt qu'une décision arbitraire.
Le chargement se fait depuis l'archive ZIP, sans exécuter de code distant.

Après déduplication, l'expérience de référence contient 3\,448 phrases.
La partition stratifiée conserve les proportions de classes dans les deux
sous-ensembles : 2\,759 exemples pour l'entraînement et 689 pour le test.
La graine pseudo-aléatoire vaut 42. Le fichier de résultats enregistre
l'empreinte SHA-256 du corpus, la configuration, les versions de Python et
des bibliothèques, ainsi que les métriques obtenues.

\subsection{Partitions et prévention de fuite}
Dans le split stratifié, les exemples sont répartis aléatoirement à l'intérieur
de chaque classe afin que les catégories minoritaires soient représentées
dans le train et le test. La graine fixe rend cette répartition reproductible,
mais elle ne transforme pas le test en évaluation temporelle. Le benchmark
propose également un split chronologique pour les fichiers CSV disposant
d'horodatages ISO-8601 : les dates les plus récentes sont réservées au test,
et les observations d'une même journée ne sont pas séparées.

La présence de phrases répétées, de formulations très proches ou d'articles
liés peut conduire à une estimation trop optimiste si des contenus liés
apparaissent dans les deux partitions. La déduplication exacte aide, sans
résoudre la fuite par quasi-doublons ni simuler le décalage temporel du marché.
Une évaluation destinée à une utilisation future devrait privilégier un test
chronologique et documenter la date de coupure.

\subsection{Métriques}
L'accuracy mesure la proportion totale de prédictions exactes. Elle est
facile à lire, mais devient trompeuse si une classe domine. Le F1 d'une classe
combine précision et rappel :
\[
  F_1 = 2\,\frac{\mathrm{précision}\times\mathrm{rappel}}
  {\mathrm{précision}+\mathrm{rappel}}.
\]
Le F1 macro est la moyenne non pondérée des F1 des trois classes. Il rend
visible la faible performance sur une classe minoritaire qui pourrait être
cachée par un bon score global. Le rapport expose aussi la précision, le
rappel, le F1 et le support de chaque classe, ainsi que la matrice de
confusion.

Les coûts sont mesurés dans le même script : temps d'entraînement ou de
chargement, temps d'inférence normalisé à 1\,000 textes, taille d'un modèle
sérialisé pour les classifieurs légers ou taille des poids pour FinBERT, et
pic d'allocations Python suivi par \texttt{tracemalloc}. Ces quantités ne
sont pas toutes équivalentes ; leurs limites sont détaillées dans la section
des résultats.

\section{Méthodes implémentées}

\subsection{Classifieur entropique Bernoulli}
Le classifieur du fichier \texttt{entropy\_sentiment.py} utilise des tokens en
minuscules et une représentation binaire par document : un mot est présent
ou absent, quelle que soit sa fréquence dans la phrase. Quelques marqueurs de
négation sont reconnus et les mots dans une portée courte après ces marqueurs
sont préfixés, par exemple \texttt{not\_good}. Il s'agit d'une heuristique
simple et non d'une analyse syntaxique de la négation.

Pour une caractéristique $t$, le gain d'information compare l'entropie des
étiquettes avant et après observation de sa présence :
\[
 IG(t)=H(Y)-P(t)H(Y\mid t)-P(\neg t)H(Y\mid\neg t).
\]
Les caractéristiques dont le gain est nul ne sont pas retenues. Les autres
sont triées par gain décroissant, avec le mot comme critère de départage
déterministe, et au plus 5\,000 sont conservées dans le benchmark.
La prise en compte de l'absence est adaptée au modèle Bernoulli : ne pas voir
un terme peut également contribuer à la prédiction.

Pour chaque classe, le modèle calcule une probabilité a posteriori lissée de
présence de chaque terme. Le score de classe combine le logarithme de l'a
priori de cette classe et les log-vraisemblances de présence ou d'absence,
pondérées par le gain d'information du terme. Cette pondération favorise les
caractéristiques informatives. Elle ne transforme toutefois pas ce modèle en
optimiseur MDL complet : il n'évalue pas explicitement un coût total de
description du modèle et des résidus.

\subsection{Classifieur par compression et approximation MDL}
La classe de compression, \texttt{CompressionMDLClassifier}, apprend un
dictionnaire par étiquette à partir des textes d'entraînement uniquement.
Les textes d'une classe sont concaténés et le dictionnaire est limité aux
32\,768 derniers octets. Pour chaque classe $c$ et un texte de test $x$, le
score utilisé est :
\[
 S(c,x) =
 \frac{|\mathrm{zlib}(D_c)|}{N_c}
 + |\mathrm{zlib}(x\mid D_c)|
 + \frac{-\log_2(N_c/N)}{8},
\]
où $D_c$ est le dictionnaire de la classe, $N_c$ le nombre de textes
d'entraînement de cette classe et $N$ le nombre total de textes
d'entraînement. La prédiction est la classe dont le score est le plus faible.
Le premier terme amortit le coût compressé du dictionnaire par les exemples
de classe ; le second mesure le résidu du texte encodé avec ce dictionnaire ;
le troisième pénalise les classes moins probables.

Cette construction traduit l'intuition MDL : réutiliser une régularité de
classe peut rendre le texte moins coûteux à encoder. C'est une approximation
expérimentale. Le choix de concaténation, la limite de taille, l'algorithme
zlib et l'amortissement par classe influencent le score. Il ne s'agit pas
d'une longueur de description universelle ni d'une garantie théorique
d'optimalité. Les dictionnaires sont construits sans utiliser le test afin
d'éviter une fuite directe des exemples évalués.

\subsection{Baseline et FinBERT}
La baseline majorité prédit pour toute phrase l'étiquette la plus fréquente
dans le train. Elle donne un seuil simple, mais sa réussite ne signifie pas
qu'elle distingue le sentiment : dans cette expérience elle ne prédit que
\emph{neutral}.

FinBERT, identifié comme \texttt{ProsusAI/finbert}, sert de référence
pré-entraînée. Il est utilisé sans fine-tuning sur Financial PhraseBank,
sur CPU, avec un maximum de 128 tokens et des lots de 16 textes. Cette
comparaison met en regard des systèmes ayant des historiques d'apprentissage
différents : FinBERT a déjà acquis de l'information pendant son
pré-entraînement, alors que les modèles légers apprennent sur la partition
d'entraînement du benchmark. Les scores ne doivent donc pas être interprétés
comme une comparaison contrôlée d'architectures recevant exactement les mêmes
données d'apprentissage.

\section{Résultats}

\subsection{Qualité de classification}
Le tableau suivant présente les scores du test de 689 phrases.

\begin{table}[ht]
\centering
\caption{Résultats de classification sur Financial PhraseBank.}
\label{tab:resultats}
\begin{tabular}{lrrr}
\toprule
Modèle & Accuracy & F1 macro & Rappel négatif\\
\midrule
Baseline majorité & 0,621 & 0,255 & 0,000\\
Entropie pondérée & 0,647 & 0,321 & 0,000\\
Compression zlib/MDL & 0,724 & 0,509 & 0,095\\
FinBERT & 0,939 & 0,931 & 0,988\\
\bottomrule
\end{tabular}
\end{table}

La baseline atteint 62,1\,\% d'accuracy parce que 428 des 689 phrases de test
sont neutres. Elle attribue les 689 exemples à cette classe et son F1 macro
reste à 25,5\,\%. Le classifieur entropique apporte un gain modeste en
accuracy, mais ne reconnaît aucun exemple négatif et n'en récupère que 18
sur 177 parmi les positifs.

Le classifieur par compression est meilleur que la baseline et le modèle
entropique sur les deux métriques globales. Son rappel positif est de 41,2\,\%
(73 exemples sur 177) et son rappel négatif de 9,5\,\% (8 sur 84). Il prédit
donc beaucoup de phrases négatives comme neutres ou positives. L'amélioration
du F1 macro à 50,9\,\% ne rend pas le modèle équilibré.

FinBERT obtient le meilleur résultat sur chaque mesure globale et un rappel
élevé dans chaque classe : 98,8\,\% pour les négatifs, 91,4\,\% pour les
neutres et 97,7\,\% pour les positifs. Dans le test étudié, il est donc le
choix le plus performant lorsque la qualité de classification prime.

\subsection{Matrice de confusion de zlib}
La matrice montre plus précisément où la méthode compacte échoue.

\begin{table}[ht]
\centering
\caption{Prédictions zlib par classe vraie (lignes) et prédite (colonnes).}
\label{tab:confusion}
\begin{tabular}{lrrr}
\toprule
Classe vraie & Prédite négative & Prédite neutre & Prédite positive\\
\midrule
Négative (84) & 8 & 58 & 18\\
Neutre (428) & 1 & 418 & 9\\
Positive (177) & 3 & 101 & 73\\
\bottomrule
\end{tabular}
\end{table}

La plupart des erreurs des classes négative et positive sont des prédictions
neutres. Cette observation suggère que la séparation par dictionnaires de
classe et l'usage des textes courts capturent assez bien les régularités
neutres, mais pas les signaux discriminants nécessaires pour les catégories
minoritaires. Le tableau n'indique pas à lui seul la cause : cela peut tenir
aux propriétés des données, à la représentation, au score zlib, ou à leur
combinaison. Il faudrait une analyse d'erreurs qualitative avant de choisir
une modification de modèle.

\subsection{Qualité, vitesse et taille}
\begin{table}[ht]
\centering
\caption{Coûts rapportés par le benchmark pilote.}
\label{tab:couts}
\begin{tabular}{lrrr}
\toprule
Modèle & Inférence (s/1\,000) & Taille mesurée (Mo) & Pic Python (Mo)\\
\midrule
Entropie pondérée & 6,43 & 0,264 & 8,90\\
Compression zlib/MDL & 3,73 & 0,099 & 0,72\\
FinBERT & 71,76 & 437,94 & 11,53\\
\bottomrule
\end{tabular}
\end{table}

Sur cette machine et pour cette exécution, zlib traite 1\,000 textes environ
19 fois plus vite que FinBERT en inférence CPU. Son état sérialisé fait
environ 0,099 Mo, contre 437,94 Mo pour les poids de FinBERT seuls. Cette
différence de taille est indicative, mais pas une comparaison strictement
homogène : d'un côté, un objet léger sérialisé ; de l'autre, les paramètres
du modèle, sans tokenizer et sans environnement d'exécution. Le rapport
signale cette différence explicitement.

De même, le pic mémoire de \texttt{tracemalloc} ne mesure que les
allocations Python observées et ne couvre pas les allocations natives de
PyTorch. Il ne représente ni la RAM maximale du processus FinBERT ni sa
mémoire totale d'inférence. Les temps dépendent du processeur, des versions
logicielles, du chargement et des conditions d'exécution. Ces nombres
caractérisent un pilote, pas un engagement de performance en production.

\section{Interprétation à la lumière du MDL}

Le résultat zlib donne un exemple concret de compromis entre modèle et
données résiduelles. Une classe fournit un dictionnaire qui aide à compresser
certaines phrases ; un texte qui partage des motifs avec ce dictionnaire peut
avoir un coût résiduel inférieur. Le coût amorti du dictionnaire et l'a priori
de classe évitent de considérer uniquement la compression du texte isolé.
La décision revient à comparer les longueurs de description approximatives
des hypothèses concurrentes.

Mais le fait d'employer un compresseur ne garantit pas, à lui seul, qu'un
système réalise l'idéal MDL. Le score ici dépend du dictionnaire construit par
concaténation, d'une fenêtre de taille fixe et d'une quantité de coût
amortie. Le modèle ne recherche pas parmi toutes les explications possibles,
ne calcule pas la complexité de Kolmogorov et n'optimise pas un code
universel. Le terme « MDL » doit donc être compris comme une inspiration et
une approximation calculable, comme l'indiquent la documentation et les
rapports du dépôt.

Le modèle entropique reflète une autre idée informationnelle : il favorise les
mots qui réduisent l'incertitude sur l'étiquette. Un gain d'information nul
signifie qu'observer la présence ou l'absence d'un terme n'améliore pas la
connaissance de la classe dans le train. La pondération rend cette intuition
opérationnelle ; elle n'encode cependant pas explicitement les données
résiduelles comme le fait une formulation MDL.

Enfin, « compréhension est compression » est une intuition utile, non une
équivalence empirique immédiate entre taille et compréhension. Une phrase
financière peut être prédite par une régularité de classe sans que le modèle
comprenne toutes ses implications. La qualité obtenue doit être mesurée sur
des exemples indépendants, et les coûts doivent être comparés selon des
définitions communes.

\section{Limites et validité}

\subsection{Limites des données}
Financial PhraseBank est un corpus annoté de phrases financières, pas un flux
continu de nouvelles de marché. Les phrases sont courtes et la tâche à trois
classes ne couvre pas toutes les distinctions pertinentes pour une décision
financière. Le split stratifié aléatoire permet une comparaison pratique,
mais ne répond pas à la question de la généralisation à une période future.
Un seul split, même reproductible, ne mesure pas la variabilité selon la
graine ou la date de coupure.

\subsection{Limites de la comparaison}
FinBERT est pré-entraîné, tandis que les méthodes légères sont entraînées sur
2\,759 exemples de train. Les méthodes ne partent pas du même volume de
connaissances. Le benchmark compare ici des systèmes utilisables en
inférence, et non des architectures soumises à des budgets d'apprentissage
identiques. Une étude complémentaire pourrait comparer les modèles légers à
une baseline TF-IDF/logistic regression et explorer des versions de FinBERT
compressées ou quantifiées.

Les mesures de taille ne comprennent pas les mêmes composants ; le pic
Python ne couvre pas toute la mémoire ; l'inférence dépend du matériel. Pour
une comparaison opérationnelle, il faudrait mesurer le processus complet
(poids, tokenizer, runtime et mémoire de pointe), dans un environnement
identique, sur les mêmes lots et avec une procédure de chauffe explicite.

\subsection{Limites des modèles}
La méthode entropique utilise une portée de négation courte et des
caractéristiques lexicales binaires. Elle ne modélise pas les dépendances
syntaxiques, les longues portées, le contexte discursif ou les termes
financiers composés de manière sophistiquée. Le modèle zlib est sensible à la
représentation textuelle et à la composition de ses dictionnaires ; il n'a pas
de mécanisme explicite de calibration et son faible rappel négatif est une
limite importante pour toute application où les pertes liées à cette classe
sont élevées.

Ces limites empêchent de conclure que la compression est suffisante pour
remplacer FinBERT sur des décisions sensibles. Un système léger pourrait
plutôt servir de filtre, de baseline, de solution embarquée sous contraintes
ou de première étape d'un pipeline, à condition d'être validé pour cet usage.

\section{Reproductibilité et structure du livrable}

Les scripts et résultats sont disponibles dans le dépôt
\href{https://github.com/EtiennePellissier/Compressed-representation}{Compressed-representation}.
Le benchmark commun se trouve dans
\texttt{Traitement/evaluation/benchmark.py}, le classifieur entropique dans
\texttt{entropy\_sentiment.py}, et les artefacts de l'expérience pilote dans
\texttt{results/financial\_phrasebank.json} et
\texttt{results/financial\_phrasebank.md}. Les tests automatisés sont dans
\texttt{tests/}.

La commande documentée pour reconstruire l'expérience télécharge l'archive
Financial PhraseBank puis passe explicitement les paramètres de corpus,
d'accord, de graine et de sortie :
\begin{lstlisting}
python Traitement/evaluation/benchmark.py \
  --phrasebank-zip /tmp/FinancialPhraseBank-v1.0.zip \
  --agreement 75 --seed 42 \
  --out results/financial_phrasebank.json
\end{lstlisting}
Le premier lancement de FinBERT nécessite également ses dépendances et le
téléchargement du modèle pré-entraîné. Les instructions exactes, notamment
l'installation CPU de PyTorch et la licence du corpus, figurent dans le
\texttt{README.md} du dépôt. Le corpus Financial PhraseBank est distribué sous
licence CC BY-NC-SA 3.0 ; toute redistribution et réutilisation doivent en
respecter les conditions et citer la publication associée.

Le chargeur lit l'archive localement et refuse les entrées mal formées ou les
étiquettes contradictoires. Les rapports contiennent l'empreinte SHA-256,
la graine, les tailles des partitions, les paramètres, les métriques et les
limites d'interprétation. Les tests de la bibliothèque standard peuvent être
lancés depuis la racine du dépôt :
\begin{lstlisting}
python -m unittest discover -s tests -v
\end{lstlisting}
La reproductibilité exacte des durées et de certains résultats de modèles
peut varier avec le matériel et les versions logicielles ; les métadonnées
enregistrées permettent de rendre cette variation visible.

\section{Perspectives}

Les prochaines étapes doivent d'abord renforcer la validité expérimentale :
\begin{enumerate}
  \item répéter le benchmark sur plusieurs graines, et rapporter la moyenne,
  la dispersion et les intervalles d'incertitude ;
  \item produire un test chronologique avec des articles dédupliqués et une
  date de coupure explicite ;
  \item ajouter une baseline lexicale TF-IDF avec régression logistique ;
  \item examiner manuellement les erreurs négatives et positives, en séparant
  les erreurs d'annotation, les ambiguïtés et les limites du modèle ;
  \item harmoniser les mesures de stockage et de mémoire, puis mesurer les
  coûts sur le matériel cible.
\end{enumerate}

Côté modélisation, il serait possible d'expérimenter d'autres normalisations,
des expressions multi-mots, des scores mieux calibrés et différents budgets
de dictionnaire. Toute modification devrait être choisie à partir du jeu
d'entraînement ou d'une validation, et non optimisée sur le test. Il faut
également vérifier si un gain sur la classe négative se maintient dans une
évaluation temporelle.

La piste conceptuelle du cours pourrait être prolongée en étudiant si des
coûts résiduels élevés signalent des textes hors domaine ou atypiques. Cela
constituerait une nouvelle tâche de détection d'anomalies, à définir et
évaluer séparément ; ce n'est pas un résultat établi par l'expérience
actuelle. De même, les méthodes de résolution d'analogies du cours peuvent
motiver la recherche de règles partageables, mais elles nécessiteraient une
implémentation et un protocole d'évaluation propres.

\section{Conclusion}

Ce travail a transformé une intuition du cours -- choisir des explications
qui décrivent les données de façon économique -- en deux approches
calculables pour le sentiment financier. Le classifieur entropique utilise
l'information mutuelle entre mots et étiquettes pour réduire et pondérer sa
représentation ; le classifieur zlib compare des coûts de compression par
classe comme approximation pratique d'une décision MDL. Un benchmark commun
les met en regard d'une baseline majorité et de FinBERT sur un corpus annoté.

Dans l'expérience pilote, zlib atteint un meilleur équilibre qualité-coût
que la méthode entropique, avec un F1 macro de 50,9\,\%, un état sérialisé
compact et un temps d'inférence CPU mesuré plus court que celui de FinBERT.
Mais son rappel négatif de 9,5\,\% révèle une faiblesse majeure, et FinBERT
reste largement supérieur en qualité avec 93,1\,\% de F1 macro. La conclusion
n'est donc pas que la compression remplace les grands modèles, mais qu'elle
offre une représentation et une famille de compromis à examiner lorsque les
ressources sont contraintes.

La suite logique est de consolider les résultats sur plusieurs partitions,
de préférer une évaluation chronologique pour les textes financiers et de
mesurer les coûts sur une base homogène. C'est seulement après ces étapes
qu'il sera possible de déterminer si une version compacte est adaptée à un
usage réel.

\section*{Références}
\addcontentsline{toc}{section}{Références}
\begin{thebibliography}{9}
\bibitem{malo}
P. Malo et al. (2014).
\newblock Good debt or bad debt: Detecting semantic orientations in economic
texts.
\newblock \emph{Journal of the Association for Information Science and
Technology}, 65.

\bibitem{araci}
D. Araci (2019).
\newblock FinBERT: Financial Sentiment Analysis with Pre-trained Language
Models.
\newblock \url{https://arxiv.org/abs/1908.10063}

\bibitem{rissanen}
J. Rissanen (1978).
\newblock Modeling by shortest data description.
\newblock \emph{Automatica}, 14(5), 465--471.

\bibitem{course}
\emph{Machine Learning and Algorithmic Information}, chapitre 4,
« Detailed outline: induction, analogies, Minimum Description Length,
prototype-based models and compression ». Support de cours.

\bibitem{phrasebank}
Financial PhraseBank, version annotée avec accord de 75\,\%.
\newblock Licence CC BY-NC-SA 3.0. Voir les références et consignes de
téléchargement dans le README du projet.
\end{thebibliography}

\end{document}
